\question[4]
Ein Von-Neumann-Rechner arbeitet in einem ständigen Befehlszyklus.
\begin{enumerate}[label=\alph*)]
\item Beschreibe die drei Phasen des Befehlszyklus.
\item Gib die Mikrobefehle an, mit denen Johnny in der ersten Phase einen Befehl holt,
und erkläre jeweils, was passiert.
\end{enumerate}

\begin{solutionbox}{8cm}
\begin{enumerate}[label=\alph*)]
\item
\begin{itemize}
\item \textbf{Holen} (Fetch): Der Befehl, auf den der Befehlszähler zeigt, wird aus dem Speicher
ins Befehlsregister geholt.
\item \textbf{Entschlüsseln} (Decode): Das Steuerwerk stellt fest, um welchen Befehl es sich handelt.
\item \textbf{Ausführen} (Execute): Der Befehl wird ausgeführt. Der Befehlszähler wird auf den
nächsten Befehl gesetzt (bei einem Sprung auf die Sprungadresse). Danach beginnt der Zyklus von vorn.
\end{itemize}
\item
\begin{verbatim}
pc->ab    Adresse des nächsten Befehls auf den Adressbus
ram->db   Befehl aus dem RAM auf den Datenbus
db->ins   Befehl ins Befehlsregister
ins->mc   Sprung zum Mikroprogramm des Befehls (Entschlüsseln)
\end{verbatim}
\end{enumerate}
\end{solutionbox}
